\section{Discussion}
\label{sec:discussion}

\subsection{Interpretation}

Our results are a caution against the premise that adaptive, per-instance correction is
worth its complexity, and a positive result about ensembling that the field already
half-knows but has not connected to agentic self-correction.

\paragraph{The benefit lives in the predictive distribution.}
Agentic forecasting systems are evaluated almost exclusively on point error. If the
effect of correction is concentrated in the predictive distribution instead of its
centre -- as our coverage results show -- a point-error benchmark is close to blind to
it. This may explain why practitioners report that reflection helps while controlled
point-accuracy comparisons show little: both observations can be correct if the
mechanism is calibration repair rather than point-estimate improvement.

\paragraph{Adaptivity did not earn its complexity here.}
We built reliability-gated correction specifically to avoid the failure modes of naive
gating: a 2\% acceptance margin that fires near-constantly, ratio-scale gains that
explode near zero, and no accounting for per-action transfer. Each fix measurably
improved the policy over its predecessor (Section~\ref{sec:pathologies}). Yet the fully
corrected policy is still dominated by the simplest static baseline in the action space:
apply the backtest-weighted ensemble unconditionally. This is a classical result in the
forecasting literature -- combination is often the single most reliable intervention
\citep{fforma2020} -- appearing here in an unfamiliar setting: it beats not just
individual models but an agent explicitly built to decide, instance by instance, when
and how to intervene. We did not find a version of adaptive gating that closes this gap,
and we report that search as unsuccessful rather than omit it.

\paragraph{The language model's problem is the decision to act, not the reflection.}
On the primary family every LLM arm is worse than never correcting
(Section~\ref{sec:llmarms}). We expected verbal self-critique to be the culprit and it
was not: the critique and single-pass arms perform within a point of each other, because
the critique barely changes what the model decides. What does the damage is a
disposition to intervene, and it has a measurable source. The model's stated risk that
a forecast is unreliable averages \calLlmMeanRisk{}, far above the rate at which large
errors actually occur. It is not indifferent to reliability; it is miscalibrated in the
direction that manufactures work.

Neither deliberation nor demonstration fixed this. The reasoning model intervened most
often of the three families, and worked examples made the critique gate worse. The
reasoning model may choose repairs better than a coin once it has decided to act, but
that is a direction, not an effect, at its sample size.

The literature's implicit theory is that agents correct badly because their critique is
unreliable, and that better reflection will fix it. Our evidence points elsewhere. With no
realizable per-series headroom (Section~\ref{sec:decomposition}), a policy can only
gain by declining to act when acting would hurt, and that is precisely what a
miscalibrated model does not do.

\paragraph{This narrows, rather than closes, the case for agentic correction.}
Our result does not show that per-instance adaptivity is worthless in general -- only
that decision rules of the kind current practice would produce (calibrated thresholds
over diagnostic signals) do not beat a strong static default in this benchmark. The gap
between ``ensemble always'' and the oracle (Figure~\ref{fig:oracle}, realizable column)
is not zero, so some instance-specific information remains unused. What we can say is
that any future adaptive policy should be benchmarked against \emph{this} baseline, not
against never-correcting, because never-correcting is the easy comparison to beat.

\paragraph{Why adaptivity loses: the gains do not transfer.}
An earlier version of this paper attributed the failure of diagnose-then-repair to two
named losses. That decomposition was built from test-selected minima and was almost
entirely winner's curse (Section~\ref{sec:decomposition}); we withdrew it. What
replaces it is simpler. The action that works best on a series at one set of origins
is a poor guide to what works best at the next, so a per-series choice made on
validation carries little information about test, and a single static action applied
everywhere does better than choosing. With four validation origins per series, the
estimate of which repair suits a series is too noisy to act on. That is the same
variance argument that explains the forecast-combination puzzle, applied to repair
selection instead of model weights.

The ground-truth diagnosis result is the part that survives unchanged and needs no
oracle. Given the true failure mode, applying its matched repair does not beat
accepting, and on CRPS it is significantly worse. Accurate diagnosis is not what is
missing.

\subsection{The measurement pathologies deserve attention on their own}

Section~\ref{sec:pathologies} reports three defects that each, independently, reversed or
concealed our conclusions before we diagnosed them. We emphasise them because they are
not specific to this study.

Any evaluation that scores a corrective action by a relative error ratio will divide by
near-zero baselines and destroy the signal it is trying to measure; the log ratio is the
appropriate scale, and the difference here is between correlations of $0.006$ and $0.43$.
Any evaluation that aggregates a heavy-tailed forecast error by the arithmetic mean is
reporting the tail. And any evaluation that quotes an oracle defined as the minimum over
many candidate outcomes is quoting a number inflated by selection; reporting a
selection-free comparator alongside it costs nothing and prevents a systematic
overstatement of what adaptive methods leave on the table.

\subsection{Relation to prior work}

Our agent is not architecturally novel, and we do not present it as such: its
diagnose--select--forecast structure follows existing agentic forecasting systems
\citep{tsscientist2025,nexus2026,castflow2026}. The difference is in the correction stage
and, more importantly, in what we measure. Where those systems reflect unconditionally or
by verbal judgement and report point error, we gate on a calibrated estimate and report
both point and probabilistic accuracy. Section~\ref{sec:decomposition}'s regret
decomposition is aimed squarely at this shared architecture: to our knowledge, no prior
agentic forecasting system has been evaluated against a ground-truth-diagnosis version of
its own correction logic, which is the only way to separate diagnosis error from a
structural weakness in diagnose-then-repair itself. Our finding -- that the pattern loses value even under perfect diagnosis, for two
separately quantified reasons -- is demonstrated for our own action space and
diagnosis-to-repair mapping, and the exact magnitudes are specific to them, but the
architectural pattern we test (a fixed mapping from diagnosis to a small candidate list,
applied without further validation) is shared broadly enough across
\citep{tsscientist2025,nexus2026,castflow2026} that the method, if not the numbers,
transfers: any such system can be probed the same way by holding its diagnosis fixed at
ground truth and asking whether its repair-selection logic recovers the headroom that
diagnosis alone made available. We can run this experiment at all only because the
counterfactual outcome tensor lets us vary the repair-selection logic while holding
diagnosis fixed at ground truth, which no existing benchmark's data supports.

\paragraph{Forecast combination, and why our headline is a rediscovery with a new
instrument.}
The result that a static combination is not beaten by adaptive per-instance selection
belongs to a literature that predates agents by half a century. \citet{bates1969}
established that combining beats choosing; \citet{clemen1989} and \citet{timmermann2006}
surveyed why, and \citet{wang2023review} brought it up to date. The forecast combination
puzzle \citep{stockwatson2004,smithwallis2009,claeskens2016} sharpens the point into
something that reads almost as a prediction of our findings: simple, equally-weighted
combinations beat optimally-weighted ones estimated from data, because the estimation
error in the weights exceeds the benefit of weighting them well. Our adaptive policies
are, structurally, weight estimators, they decide per instance which repair to trust, and they lose to an unweighted combination in exactly the way that literature says
they should. \citet{hibon2005} asked our question directly in the pre-agentic setting and
reached the same answer.

We are therefore not claiming to have discovered that combination wins. What we claim is
that the agentic forecasting literature has been constructing increasingly elaborate
adaptive correction machinery without measuring it against this prior, and that the
counterfactual tensor supplies something the classical literature could not: it holds
outcomes for corrective \emph{actions}, not just for alternative forecasts, so the
shortfall can be decomposed into named architectural causes rather than merely observed
as an aggregate loss. The combination literature says adaptive selection tends to lose.
It does not say \emph{which part} of a diagnose-then-repair pipeline loses it, or how
much each part costs. Section~\ref{sec:decomposition} does.

The reliability estimator descends from feature-based forecast performance prediction
\citep{fformpp2019} and forecast stress testing \citep{mast2025}, which established that
failure is predictable from features. We close the loop by using that prediction to
control an agent, and we find that the prediction is accurate while the resulting
intervention is worth making far less often than the literature assumes.

Relative to existing time-series agent benchmarks
\citep{temporalbench2026,timesagemt2026,timeseriesgym2025,aion2026}, TS-AgentBench
evaluates a different object: the control decision, which requires counterfactual
outcomes the others do not record. The closest work in spirit is
\citet{agentabstain2026}, which asks whether tool-using agents know when not to act; that
work addresses general agentic settings without counterfactual outcome data, whereas here
the counterfactual is computed exactly.

\subsection{Limitations}

\paragraph{Computational scale.}
All experiments ran on a CPU-only machine. Neural forecasters are small and trained on a
subsample, and the foundation-model baseline is optional. We therefore make no claim
about the best available forecasters, only about correction policies applied to a fair,
identically-treated pool.

\paragraph{The action space is finite and hand-designed.}
Ten typed actions is what makes the counterfactual tensor computable, and it is also a
restriction. A repair an expert would make but which is absent from the space is
invisible to our oracle, so headroom estimates are relative to this space.

\paragraph{The result does not depend on the base forecast having been adaptively
chosen.}
Our base forecast is chosen per instance by rolling-origin backtest over a nine-model
pool. That is itself a strong adaptive procedure, and it was the largest threat to
every claim above: it plausibly leaves correction little to repair, which would
confound ``correction does not work'' with ``correction had nothing to fix''. A deployed
single-model pipeline is a materially easier target, and correction might have paid
there.

We tested this with a second condition that fixes the base forecaster to ETS, the model
the backtest selector picks most often and a realistic default for a deployed pipeline,
with every other factor held identical (\texttt{configs/fixed\_base.yaml}). The two
runs differ in exactly one scientific setting; this was verified mechanically rather
than asserted, by confirming identical instance sets and identical history-derived
diagnostics before any comparison was reported. Both readings of the outcome were fixed
in the analysis plan before the run completed (deviation~D11).

Table~\ref{tab:fixedbase} gives the answer. Under a fixed ETS base no policy
significantly improves on never correcting, exactly as under the adaptive base.
Unconditional correction remains significantly harmful (\fbAlwaysCorrectFixedPct\%,
[\fbAlwaysCorrectFixedCIlo, \fbAlwaysCorrectFixedCIhi]\%, against
\fbAlwaysCorrectAdaptivePct\% on the adaptive base), and the calibrated gate again fails
to separate from zero (\fbRgcFixedPct\%, [\fbRgcFixedCIlo, \fbRgcFixedCIhi]\%). The
hypothesis that correction failed only because the adaptive base left nothing to fix is
rejected. The conclusions above hold for a single-model pipeline, not only for the harder
backtest-selected case.

Each column of the table is a difference from never correcting \emph{within its own
run}. The two runs have different base forecasts and therefore different references, so
their levels are never compared, only their differences set side by side.

\begin{table}[htbp]
\centering
\caption{Correction under an adaptive, backtest-selected base beside the same policies
under a fixed ETS base, on identical series, splits and seeds. Each cell is the
series-level MASE difference from never correcting within its own run, with a 95\%
bootstrap interval over series; $\dagger$ marks an interval excluding zero. Positive is
worse. Levels are not shown because the two runs have different references.}
\label{tab:fixedbase}
\small
\fittable{\begin{tabular}{lrr}
\toprule
Policy & Adaptive (backtest) base & Fixed ETS base \\
\midrule
always correct & 8.13\%$^\dagger$~[2.73, 15.99] & 10.59\%$^\dagger$~[1.21, 27.56] \\
best fixed & 8.16\%$^\dagger$~[2.76, 16.05] & 10.52\%$^\dagger$~[1.15, 27.49] \\
threshold single & 5.48\%$^\dagger$~[1.46, 10.88] & 3.95\%$^\dagger$~[0.53, 9.29] \\
random correct & 2.72\%$^\dagger$~[2.06, 3.46] & 3.59\%$^\dagger$~[2.88, 4.40] \\
rgc & 0.14\%~[-1.84, 3.03] & -0.49\%~[-1.64, 0.67] \\
\bottomrule
\end{tabular}
}
\end{table}

\paragraph{Post-hoc analyses.}
The calibrated policy, the log-ratio gain scale, the robust aggregation and the
fixed-base condition were all introduced after the first test evaluation. Their
hyperparameters were selected on validation and each was evaluated on test once, but they
are exploratory rather than confirmatory and are labelled so throughout. The deviations
are itemised in the released analysis plan.

\paragraph{Ground truth for failure modes exists only on synthetic data.}
Real series carry no known change points or outlier labels, so failure-mode labels there
are derived from the same detectors the agent uses. Causal claims rest on synthetic data.

\paragraph{LLM scope.}
The LLM arms cover three model families --- a small hosted proprietary model, an
open-weight model of different lineage, and an open-weight reasoning model at medium
effort --- at temperature 0, under free-tier API budgets that cap each family to one
instance per series over part of the benchmark. The statistical weight sits with the
first, at \nseries{} series; the other two cover 140 and \reasonSeries{} series and are
powered to replicate a direction, not to establish an effect alone. Three families rule
out a single weak model, an absence of deliberation, and the zero-shot setting as
explanations; that is not enough to characterise language models in general, and we do not
claim it is.

What remains untested is a frontier proprietary reasoning model, multi-turn deliberation,
and tool access during the decision. We also did not engineer the prompt against this
benchmark: the evidence block is a fixed template, and the eight-example condition is a
control for supervision, not an attempt at optimisation. A system built around a prompt
tuned on these diagnostics might do better. The arm originally planned for the reasoning
family, \texttt{gpt-oss-120b}, exhausted its daily token allowance mid-run and was
replaced by the 20B model of the same family at the same effort setting; the substitution
is recorded in deviation~D9 and the larger model remains untested.

\paragraph{Scale.}
\nseries{} series is modest next to the large forecast-evaluation suites
\citep{gifteval2024,monash2021}, and the reason is the counterfactual tensor itself:
scoring every action on every instance costs roughly an order of magnitude more compute
than scoring one, which is the price of being able to ask decision-level questions at
all. The series-level intervals we report already account for this; several of them
include zero, and we say so instead of reaching for the instance-level test that would
not have. Where an interval is wide, this benchmark cannot resolve the effect; that is not evidence the effect is absent.

\subsection{Future work}

The regret decomposition points at two concrete, separately testable fixes instead of a
vague call for "better adaptivity." First, replace the fixed canonical mapping with
validated selection among a failure mode's candidate repairs -- our calibrated policy
already does this and is measurably less harmful than the idealised canonical-fix policy
as a result, though it still under-performs the static ensemble. Second, and
un-addressed by anything we built, replace the hand-designed diagnosis-to-repair
candidate list with one learned from the counterfactual tensor itself: the benchmark
supplies exactly the (failure signature, action, outcome) tuples needed to learn which
repairs actually help under which conditions, rather than which repairs a domain expert
guesses should help. The replay property means candidate mappings can be compared without
refitting, so this is a tractable next experiment on the released artifact instead of a
new benchmark. More generally, the tensor supplies exactly the (state, action, outcome)
tuples that off-policy evaluation and offline reinforcement learning require for
correction policies of any design. Given our finding that the benefit is probabilistic,
an action space designed for distributional repair --- richer conformal schemes, adaptive
interval widening, distributional ensembling --- is likely to have more headroom than one
built around alternative point estimates.

\section{Conclusion}
\label{sec:conclusion}

Self-correction has become a standard component of forecasting agents without evidence
that it works, because measuring it requires knowing what would have happened under the
corrections not taken. We built a benchmark that records exactly that.

Its first lesson is about measurement. The oracle headroom that makes correction look
promising is almost entirely selection: chosen on test outcomes the best repair per
instance looks transformative, and chosen on earlier origins it recovers nothing. A
correction method judged by how much of a test-selected oracle gap it closes may have
closed none of the realizable one. Correction methods should be measured against a
validation-selected oracle and against a static default, not against doing nothing.

Its second lesson is that the realizable gain comes from choosing once. A single static
ensemble improves probabilistic accuracy, holds on real data alone, and is what
validation selects; choosing a repair per series does worse. That is the
forecast-combination puzzle \citep{stockwatson2004,smithwallis2009,claeskens2016}
reappearing at the level of repair selection, and its direction is what that literature
predicts \citep{bates1969,timmermann2006,wang2023review}. Granting a perfect diagnosis
does not change it: applying the matched repair still does not help.

Its third lesson concerns the systems the field is building. When a language model makes
the decision, it corrects too often, because it believes its forecasts are far less
reliable than they are, and asking it to reflect does not change that. On the primary
model family this makes forecasts measurably worse; two smaller families point the same
way. None of these conclusions depends on the base forecast having been adaptively
chosen.

The same risk estimate that fails to drive correction is a well-calibrated basis for
triage, a usable result that asks nothing of the correction machinery. We release the
benchmark and the outcome tensor so that correction policies, including ones that
disagree with these conclusions, can be evaluated against the same counterfactual
outcomes.
